%-----------------------------------------------------------------------------------------------------------------------------------------------%
%	The MIT License (MIT)
%
%	Copyright (c) 2021 Jitin Nair
%
%	Permission is hereby granted, free of charge, to any person obtaining a copy
%	of this software and associated documentation files (the "Software"), to deal
%	in the Software without restriction, including without limitation the rights
%	to use, copy, modify, merge, publish, distribute, sublicense, and/or sell
%	copies of the Software, and to permit persons to whom the Software is
%	furnished to do so, subject to the following conditions:
%	
%	THE SOFTWARE IS PROVIDED "AS IS", WITHOUT WARRANTY OF ANY KIND, EXPRESS OR
%	IMPLIED, INCLUDING BUT NOT LIMITED TO THE WARRANTIES OF MERCHANTABILITY,
%	FITNESS FOR A PARTICULAR PURPOSE AND NONINFRINGEMENT. IN NO EVENT SHALL THE
%	AUTHORS OR COPYRIGHT HOLDERS BE LIABLE FOR ANY CLAIM, DAMAGES OR OTHER
%	LIABILITY, WHETHER IN AN ACTION OF CONTRACT, TORT OR OTHERWISE, ARISING FROM,
%	OUT OF OR IN CONNECTION WITH THE SOFTWARE OR THE USE OR OTHER DEALINGS IN
%	THE SOFTWARE.
%	
%
%-----------------------------------------------------------------------------------------------------------------------------------------------%

%----------------------------------------------------------------------------------------
%	DOCUMENT DEFINITION
%----------------------------------------------------------------------------------------

% article class because we want to fully customize the page and not use a cv template
\documentclass[a4paper,5pt]{article}

%----------------------------------------------------------------------------------------
%	FONT
%----------------------------------------------------------------------------------------

% % fontspec allows you to use TTF/OTF fonts directly
% \usepackage{fontspec}
% \defaultfontfeatures{Ligatures=TeX}

% % modified for ShareLaTeX use
% \setmainfont[
% SmallCapsFont = Fontin-SmallCaps.otf,
% BoldFont = Fontin-Bold.otf,
% ItalicFont = Fontin-Italic.otf
% ]
% {Fontin.otf}

%----------------------------------------------------------------------------------------
%	PACKAGES
%----------------------------------------------------------------------------------------
\usepackage{url}
\usepackage{parskip} 	

%other packages for formatting
\RequirePackage{color}
\RequirePackage{graphicx}
\usepackage[usenames,dvipsnames]{xcolor}
\usepackage[scale=0.9]{geometry}

%tabularx environment
\usepackage{tabularx}

%for lists within experience section
\usepackage{enumitem}

% centered version of 'X' col. type
\newcolumntype{C}{>{\centering\arraybackslash}X} 

%to prevent spillover of tabular into next pages
\usepackage{supertabular}
\usepackage{tabularx}
\newlength{\fullcollw}
\setlength{\fullcollw}{0.47\textwidth}

%custom \section
\usepackage{titlesec}				
\usepackage{multicol}
\usepackage{multirow}

%CV Sections inspired by: 
%http://stefano.italians.nl/archives/26
\titleformat{\section}{\Large\scshape\raggedright}{}{0em}{}[\titlerule]
\titlespacing{\section}{0pt}{10pt}{10pt}

%for publications
\usepackage[style=authoryear,sorting=ynt, maxbibnames=2]{biblatex}

%Setup hyperref package, and colours for links
\usepackage[unicode, draft=false]{hyperref}
\definecolor{linkcolour}{rgb}{0,0.2,0.6}
\hypersetup{colorlinks,breaklinks,urlcolor=linkcolour,linkcolor=linkcolour}
\addbibresource{citations.bib}
\setlength\bibitemsep{1em}

%for social icons
\usepackage{fontawesome5}

%debug page outer frames
%\usepackage{showframe}


% job listing environments
\newenvironment{jobshort}[2]
    {
    \begin{tabularx}{\linewidth}{@{}l X r@{}}
    \textbf{#1} & \hfill &  #2 \\[3.75pt]
    \end{tabularx}
    }
    {
    }

\newenvironment{joblong}[2]
    {
    \begin{tabularx}{\linewidth}{@{}l X r@{}}
    \textbf{#1} & \hfill &  #2 \\[3.75pt]
    \end{tabularx}
    \begin{minipage}[t]{\linewidth}
    \begin{itemize}[nosep,after=\strut, leftmargin=1em, itemsep=3pt,label=--]
    }
    {
    \end{itemize}
    \end{minipage}    
    }



%----------------------------------------------------------------------------------------
%	BEGIN DOCUMENT
%----------------------------------------------------------------------------------------
\begin{document}

% non-numbered pages
\pagestyle{empty} 

%----------------------------------------------------------------------------------------
%	TITLE
%----------------------------------------------------------------------------------------

% \begin{tabularx}{\linewidth}{ @{}X X@{} }
% \huge{Your Name}\vspace{2pt} & \hfill \emoji{incoming-envelope} email@email.com \\
% \raisebox{-0.05\height}\faGithub\ username \ | \
% \raisebox{-0.00\height}\faLinkedin\ username \ | \ \raisebox{-0.05\height}\faGlobe \ mysite.com  & \hfill \emoji{calling} number
% \end{tabularx}

\begin{tabularx}{\linewidth}{@{} C @{}}
\Huge{Shirley Katherine Maldonado} \\[6pt]
{Analista de datos $|$ Cientifica de datos }\\[8pt]
%\href{https://github.com/username}{\raisebox{-0.05\height}\faGithub\ username} \ $|$ \ 
\href{https://www.linkedin.com/in/katherine-maldonado-64424b205/}{\raisebox{-0.05\height}\faLinkedin\ Shirley Maldonado} \ $|$ \ 
%\href{https://mysite.com}{\raisebox{-0.05\height}\faGlobe \ mysite.com} \ $|$ \ 
\href{mailto:kathe.meza.27@gmail.com}{\raisebox{-0.05\height}\faEnvelope \ kathe.meza.27@gmail.com} \ $|$ \ 
\href{tel:+57 3209593820}{\raisebox{-0.05\height}\faMobile \ +57 3209593820}\ $|$ \ \href{tel:+57 3209593820}{\raisebox{-0.05\height}\faMobile \ +57 3249176840} \\
\end{tabularx}

%----------------------------------------------------------------------------------------
% EXPERIENCE SECTIONS
%----------------------------------------------------------------------------------------

%Interests/ Keywords/ Summary
\section{Perfil}
Analista y científica de datos con formación en Astronomía, con experiencia práctica en el sector retail. He desarrollado modelos predictivos del NPS y análisis de reputación de marca cuyos resultados fueron adoptados por el equipo como marcos de priorización operativa. Manejo Python, R, SQL y Excel para extracción, transformación y modelado de datos, y Power BI para comunicación de resultados a audiencias de negocio. Mi enfoque combina solidez estadística con la capacidad de traducir hallazgos técnicos en decisiones accionables.
\href{https://kikikath.github.io/Hoja-de-Vida/cv.pdf}{CV}

%Experience
\section{Experiencia laboral}

\begin{joblong}{Practicante – Investigación de Mercado y Ciencia de Datos}{oct 2025 - mar 2026}

\item Desarrollé un modelo predictivo del NPS omnicanal utilizando Gradient Boosting (R² entre 0.63 y 0.79 en conjunto de prueba), identificando el cumplimiento en tiempos de entrega de producto como el principal driver de satisfacción y lealtad del cliente. Los resultados fueron presentados al equipo y adoptados como marco de priorización operativa.
\item Ejecuté la optimización de la encuesta de experiencia del cliente mediante PCA y análisis de correlaciones, identificando ítems redundantes o con baja varianza (como Medios de pago disponibles y Conocimiento del personal) cuya eliminación reduce la carga de respuesta sin pérdida de poder explicativo (varianza explicada por Momento de Verdad: 75–85%).
\item Diseñé un análisis de reputación de marca con regresión logística ordinal (en lugar del enfoque estándar por importancia de permutación) para cuantificar el impacto de cada variable sobre la confianza del consumidor. Esta elección metodológica, propuesta por mí y validada por el equipo, permitió obtener pesos directamente interpretables que se agregaron a nivel de driver y se tradujeron en un gráfico de desempeño vs. impacto coherente entre públicos y niveles de análisis.
\item Analicé datos Nielsen para benchmarking competitivo y apoyé al equipo de Ciencia de Datos en métricas de percepción de marca.
\end{joblong}
  
%Projects
\section{Proyectos}

\begin{tabularx}{\linewidth}{ @{}l r@{} }
\textbf{Análisis de Series de Tiempo y Predicción del Valor de la Energía Eléctrica} & \hfill \href{https://some-link.com}{ene 2025} \\[3.75pt]
\multicolumn{2}{@{}X@{}}{Proyecto académico en equipo en donde procesamos datos históricos de XM, construimos modelos LSTM en TensorFlow/Keras para pronosticar el valor del kWh a corto y mediano plazo, y evaluamos el desempeño con métricas MAE y RMSE.}  \\
\end{tabularx}

\begin{tabularx}{\linewidth}{ @{}l r@{} }
\textbf{Diseño de Base de Datos para Gestión del Riesgo y Desastre en Antioquia} & \hfill \href{https://some-link.com}{may 2024} \\[3.75pt]
\multicolumn{2}{@{}X@{}}{Proyecto de cierre del Bootcamp Data Analytics de Betek, presentado en el Open House de la empresa ante audiencia externa. Diseñé el modelo relacional en MySQL, realicé el análisis exploratorio con Python y construí un tablero en Power BI con el histórico de desastres y zonas de riesgo del departamento de Antioquia.}  \\
\end{tabularx}

%----------------------------------------------------------------------------------------
%	EDUCATION
%----------------------------------------------------------------------------------------
\section{Educación}
\begin{tabularx}{\linewidth}{@{}l X@{} l@{}}	
mar 2019 - mar 2026 & Pregrado en Astronomía & \textbf{Universidad de Antioquia} \hfill \normalsize %(GPA: 4.0/4.0) \\

\end{tabularx}

%----------------------------------------------------------------------------------------
%	PUBLICATIONS
%----------------------------------------------------------------------------------------
\section{Educación complementaria}
\begin{tabularx}{\linewidth}{@{}l X@{} l@{}}

oct 2024 - dic 2024 & Curso Power BI y Visualización de Datos & 
\textbf{Universidad de Antioquia} \hfill \normalsize \\

nov 2023 - may 2024 & Bootcamp Data Analytics & \textbf{Betek} \hfill \normalsize \\

nov 2023 - feb 2024 & Diplomado en Análisis de Datos y Machine Learning con Python & \textbf{Universidad de Antioquia} \hfill \normalsize \\

nov 2022 - dic 2022 & Fundamentos en Analítica de Datos  & \textbf{DS4A} \hfill \normalsize \\

\end{tabularx}


\section{Idiomas}
\begin{tabularx}{\linewidth}{@{}l X@{}}
Inglés & \normalsize{Avanzado}
\end{tabularx}


%----------------------------------------------------------------------------------------
%	SKILLS
%----------------------------------------------------------------------------------------
\section{Habilidades Técnicas}
\begin{tabularx}{\linewidth}{@{}l X@{}}
Programación:  &  \normalsize{Python, R, SQL.}\\
Machine learning:  &  \normalsize{ Histogram Gradient Boosting, Regresión logistica ordinal, Clasificación, Árboles de Decisión, LSTM.}\\  
Métodos estadísticos:  &  \normalsize{PCA, Análisis de Correspondencias Múltiples (ACM), Algoritmo PC (descubrimiento
causal).}\\
Interpretabilidad:  &  \normalsize{SHAP, importancia por permutación.}\\
Bases de Datos:  &  \normalsize{MySQL, Excel y SAS.}\\
Visualización:  &  \normalsize{Power BI, Matplotlib, Seaborn, Qualtrics.}\\
\end{tabularx}


\vfill
\center{\footnotesize Last updated: \today}

\end{document}
